\documentclass[10pt,a4paper]{article}
\usepackage[UTF8]{ctex}
\usepackage{amsmath,amssymb,geometry,xcolor,listings,tcolorbox,multicol,enumitem}
\geometry{margin=1.2cm,top=1cm,bottom=0.8cm}
\setlength{\parskip}{0pt}\setlength{\parsep}{0pt}\setlength{\itemsep}{0pt}
\setlist{nosep,leftmargin=*}

\lstset{language=C,basicstyle=\ttfamily\scriptsize,keywordstyle=\color{blue}\bfseries,
  numbers=none,frame=single,breaklines=true,showspaces=false,tabsize=2,
  aboveskip=2pt,belowskip=2pt}

\newtcolorbox{keybox}[1][]{colback=yellow!3!white,colframe=yellow!70!black,
  title=\textbf{考点},fonttitle=\bfseries\small,size=small,top=2pt,bottom=2pt,#1}
\newtcolorbox{bugbox}[1][]{colback=red!3!white,colframe=red!60!black,
  title=\textbf{易错点},fonttitle=\bfseries\small,size=small,top=2pt,bottom=2pt,#1}

\title{\textbf{二叉树 + 二叉搜索树 (BST)}}
\author{}\date{}

\begin{document}
\maketitle\vspace{-8pt}

\begin{multicols}{2}

% ====== 术语 ======
\section*{1. 树的术语}

\begin{tabular}{|l|l|}
\hline
\textbf{术语} & \textbf{定义} \\
\hline
Degree(节点的度) & 该节点的直接孩子数 \\
Degree of Tree & 所有节点度的最大值 \\
Parent/Child/Sibling & 父/子/兄弟 \\
Leaf/Leaf & 度为0的节点（没有孩子） \\
Path & 节点序列 $n_1,n_2,\dots,n_k$ \\
Path Length & 路径上的\textbf{边数}（不是节点数） \\
Depth & 从根向下的边数，根=0 \\
Height & 到最深叶子的边数，叶子=0 \\
Ancestor & 从节点到根的路径上所有节点 \\
Descendant & 该节点子树中所有节点 \\
\hline
\end{tabular}

\begin{keybox}
\textbf{深度 $\neq$ 高度：} 深度从上往下，高度从下往上。\\
\textbf{叶子 $\neq$ 孩子：} 叶子是度=0的节点，孩子是有父节点的节点。\\
\textbf{路径长度 = 边数（不是节点数）}
\end{keybox}

% ====== 二叉树性质 ======
\section*{2. 二叉树三性质}

\begin{enumerate}
\item 第 $i$ 层最多 $2^{i-1}$ 个节点 ($i\ge 1$)
\item 深度 $k$ 的树最多 $2^{k}-1$ 个节点
\item \textbf{$n_0 = n_2 + 1$}（叶子数 = 度2节点数+1）
\end{enumerate}

\textbf{性质3证明：} $n=n_0+n_1+n_2$，$B=n_1+2n_2$，$n=B+1$。代入得证。\\
\textbf{空链接数：} $n$ 个节点的二叉树有 $n+1$ 个 NULL 链接。

\subsection*{二叉树类型}

\textbf{斜二叉树：} 全左或全右 → 最坏深度 $O(N)$。\\
\textbf{完全二叉树：} 除最底层，其他全满。最底层从左到右连续。

\subsection*{表达式树}

叶子=操作数，内部节点=运算符。\\
后缀表达式 → 表达式树：用栈构建。读到操作数 push，读到运算符 pop 两个当孩子再 push。

\textbf{三种遍历对应：} 中序→中缀，后序→后缀，先序→前缀。

\columnbreak

% ====== 遍历 ======
\section*{3. 四大遍历（考试必默写）}

\begin{lstlisting}
// 先序：根→左→右
void preorder(TreeNode* r) {
    if(!r) return;
    printf("%d ",r->data);
    preorder(r->left);
    preorder(r->right);
}

// 中序：左→根→右
void inorder(TreeNode* r) {
    if(!r) return;
    inorder(r->left);
    printf("%d ",r->data);
    inorder(r->right);
}

// 后序：左→右→根
void postorder(TreeNode* r) {
    if(!r) return;
    postorder(r->left);
    postorder(r->right);
    printf("%d ",r->data);
}
\end{lstlisting}

\textbf{关键：} 三种递归只差 printf 的位置。

\subsection*{中序遍历 — 迭代版}

\begin{lstlisting}
void iterInorder(TreeNode* r) {
    Stack S; init(&S);
    for(;;){  // 也可 while(t||!isEmpty(&S))
        for(; r; r=r->left) push(&S,r);
        if(isEmpty(&S)) break;
        r=pop(&S); visit(r);
        r=r->right;
    }
}
\end{lstlisting}

\subsection*{层序遍历 — 用队列}

\begin{lstlisting}
void levelorder(TreeNode* r) {
    if(!r)return; Queue q; init(&q);
    enqueue(&q,r);
    while(!isEmpty(&q)){
        TreeNode* t=dequeue(&q);
        visit(t);
        if(t->left) enqueue(&q,t->left);
        if(t->right) enqueue(&q,t->right);
    }
}
\end{lstlisting}

\subsection*{遍历结果速记}

\begin{center}\small
\begin{tabular}{|c|c|c|}
\hline
\textbf{遍历} & \textbf{顺序} & \textbf{用途} \\
\hline
先序 (Pre) & 根→左→右 & 目录列表 \\
后序 (Post) & 左→右→根 & 目录大小、表达式求值 \\
中序 (In) & 左→根→右 & BST 输出有序序列 \\
层序 & BFS & 逐层遍历 \\
\hline
\end{tabular}
\end{center}

% ====== 线索二叉树 ======
\section*{4. 线索二叉树}

\textbf{动机：} $n$ 个节点的二叉树有 $n+1$ 个 NULL。用来指向中序前驱/后继。

\textbf{规则：}
\begin{enumerate}
\item Tree上->Left=NULL → 指向中序前驱
\item Tree上->Right=NULL → 指向中序后继
\item 加头节点，无悬空线程
\end{enumerate}

节点结构：\texttt{data, left, right, leftThread, rightThread}

% ====== BST ======
\section*{5. 二叉搜索树 (BST)}

\textbf{定义：} 左子树所有值 < 根 < 右子树所有值。键值互异（或可处理重复）。

\subsection*{Find — $O(h)$}

\begin{lstlisting}
TreeNode* find(TreeNode* r, int x) {
    if(!r) return NULL;
    if(x<r->data) return find(r->left,x);
    if(x>r->data) return find(r->right,x);
    return r;
}
// 迭代版：while(r){if(x<r->data)r=r->left;else if...else return r;}
\end{lstlisting}

\subsection*{FindMin / FindMax — 一直左 / 一直右}

\begin{lstlisting}
TreeNode* findMin(TreeNode* r) {
    if(!r)return NULL;
    if(!r->left)return r;
    return findMin(r->left);
}
TreeNode* findMax(TreeNode* r) {
    if(!r)return NULL;
    while(r->right) r=r->right;
    return r;
}
\end{lstlisting}

\subsection*{Insert — 在NULL处创建新节点}

\begin{lstlisting}
TreeNode* insert(TreeNode* r, int x) {
    if(!r) {
        TreeNode* n=(TreeNode*)malloc(sizeof(TreeNode));
        n->data=x; n->left=n->right=NULL;
        return n;
    }
    if(x<r->data) r->left=insert(r->left,x);
    else if(x>r->data) r->right=insert(r->right,x);
    return r;
}
\end{lstlisting}

\begin{bugbox}
r为NULL时直接 \texttt{r->data=x} → 段错误。必须先malloc。\\
忘了 \texttt{root->left=insert(...)} 来接返回值。
\end{bugbox}

\subsection*{Delete — 三情况}

\begin{lstlisting}
TreeNode* delete(TreeNode* r, int x) {
    if(!r)return NULL;
    if(x<r->data) r->left=delete(r->left,x);
    else if(x>r->data) r->right=delete(r->right,x);
    else { // 找到了！
        // Case 1+2: 0或1个孩子
        if(!r->left){TreeNode*t=r->right;free(r);return t;}
        if(!r->right){TreeNode*t=r->left;free(r);return t;}
        // Case 3: 两个孩子 → 找替身
        TreeNode* min=findMin(r->right);     // 右子树最小
        r->data=min->data;
        r->right=delete(r->right,min->data); // 递归删替身
    }
    return r;
}
\end{lstlisting}

\begin{keybox}
\textbf{Case 3 核心：} 找左子树最大 或 右子树最小 替换值，然后删替身（替身必是 Case 1 或 2）。\\
\textbf{不要找整棵树的最大！} 应该找左子树的最大。
\end{keybox}

\begin{bugbox}
忘了 free。\\
Case 2 的 left/right 写反（只有一个左孩子却返回了右孩子）。\\
\texttt{findMax(r)} 找整棵树最大值 → 应 \texttt{findMax(r->left)}。\\
\texttt{delete(max,max->data)} 参数错误 → 应 \texttt{r->left=delete(r->left,max->data)}。
\end{bugbox}

\subsection*{平均情况分析}

BST的深度取决于插入顺序。\\
随机插入 → 平均深度 $O(\log N)$；升序/降序插入 → 深度 $O(N)$（退化为链表）。

\subsection*{懒惰删除}

加 flag 标记删除，不真正 free。\\
好处：不用 malloc 就能重插入。坏处：dead node 太多影响效率。

\end{multicols}

\section*{二叉树操作复杂度}

{\small
\begin{tabular}{|l|c|c|}
\hline
\textbf{操作} & \textbf{时间复杂度} & \textbf{备注} \\
\hline
Find / FindMin / FindMax & $O(h)$ & h=树高 \\
Insert & $O(h)$ & 在 NULL 处创建 \\
Delete & $O(h)$ & Case 3 转 Case 1/2 \\
四种递归遍历 & $O(N)$ & 每个节点访问一次 \\
层序遍历 & $O(N)$ & 用队列 \\
BST 随机插入平均深度 & $O(\log N)$ & \\
BST 最坏深度 & $O(N)$ & 有序插入 \\
\hline
\end{tabular}
}

\end{document}
